\documentclass[12pt]{article}
\usepackage[T1]{fontenc}
\usepackage[utf8]{inputenc}
\usepackage{lmodern}
\usepackage{textcomp}
\usepackage{enumitem}
\usepackage{geometry}
\usepackage{graphicx}
\usepackage{amsmath, amssymb}
\geometry{margin=1in}
\begin{document}
\begin{center}
Economics - Graduate Level Assessment 4 \
Total Marks: 40 \
Answer all questions.
\end{center}

\section*{Multiple Choice (10 marks)}
\begin{enumerate}[label=\arabic*.]
\item Explain the two arguments that have always been advanced fora special policy towards farmers.
\begin{enumerate}[label=(\alph*)]
    \item Climate change adaptation and biotechnology
    \item International trade and subsidies competition
    \item Agricultural sustainability and urbanization pressures
    \item Land reform and rural infrastructure development
    \item Technological innovation and market expansion
\end{enumerate}
\item There is a strong demand for welders in California but Bill an unemployed welder lives in New York. Bill is
\begin{enumerate}[label=(\alph*)]
    \item frictionally unemployed.
    \item voluntarily unemployed.
    \item considered to be a hidden worker.
    \item underemployed.
    \item structurally unemployed.
\end{enumerate}
\item If nominal GDP equals \$6000 and the GDP deflator equals 200 then real GDP equals
\begin{enumerate}[label=(\alph*)]
    \item \$4,000
    \item \$6,000
    \item \$12,000
    \item \$2,000
    \item \$1,200
\end{enumerate}
\item What does the Harrod-Domar model, taken with the Keynesian theory of savings, imply for the growth rates of poor countries?
\begin{enumerate}[label=(\alph*)]
    \item Rich and poor countries will experience the same growth rates regardless of savings rates
    \item Poor countries will have fluctuating growth rates independent of savings
    \item The level of technology determines growth rates, not savings
    \item The level of income has no impact on savings
    \item Savings have no correlation with growth rates in any country
\end{enumerate}
\item Which of the following will increase wages for tuba makers?
\begin{enumerate}[label=(\alph*)]
    \item A decrease in the cost of tuba-making materials
    \item An increase in the number of graduates at tuba maker training school
    \item An increase in the tax on tubas
    \item An increase in the number of tuba maker training schools
    \item An increase in the price of tubas
\end{enumerate}
\end{enumerate}

\section*{True / False (10 marks)}
\textit{Answer True or False}
\begin{enumerate}[label=\arabic*.]
\setcounter{enumi}{5}
\item True or False: In a pure capitalism system, consumer preferences may inadequately address the production of goods and services that cannot be individually provided.
\item True or False: A celebrity's income from endorsements is classified as transfer earnings due to its dependence on market value rather than uniform pay.
\item True or False: A kinked demand curve implies that firms do not adjust output in response to changes in competitors' prices.
\item True or False: Per capita figures measure growth more accurately and reflect changes in individual welfare within an economy.
\item True or False: The trade deficit account measures only the import and export of goods and services, excluding financial transactions between nations.
\end{enumerate}

\section*{Long Form Response (20 marks)}
\textit{Answer each question with a clear and complete explanation.}
\begin{enumerate}[label=\arabic*.]
\setcounter{enumi}{10}
\item Analyze the loan transaction between Mr. Riley and Mr. Gillis, calculating both the nominal and real interest rates while discussing the implications of inflation on the value of money over the period.
\item Discuss the implications for transfer earnings and economic rent when land is designated for a single use, providing a detailed analysis of the underlying economic principles.
\end{enumerate}

\vfill
\noindent\textit{End of Paper}
\end{document}